% !TeX root = ../../Book3.tex
% 纲要标题：### 1.3 素谱的拓扑
% 纲要位置：工作记录/计划纲要.md，第 1828 行。
% 纲要细目（写作范围）：
% * Zariski 闭集与基本开集
% * 拟紧性（quasi-compactness，不附加 Hausdorff 条件）
% * 不可约闭集、一般点与闭点
% * Noether 素谱与有限不可约分量

\section{素谱的拓扑}
\label{b3:sec:0103}

把素理想作为点，还要指定哪些点属于同一个闭集。理想的包含关系给出一种自然规则：把包含同一理想的素理想收集起来。这个拓扑通常不是 Hausdorff 拓扑\footnote{豪斯多夫（Felix Hausdorff，1868-1942），德国数学家，建立一般拓扑中的分离公理等基本概念。}，甚至不同点之间可能有特殊化关系；这些现象正是需要保留的代数信息。

\subsection{Zariski 闭集与基本开集}
\label{b3:subsec:0103:01}

\begin{definition}\label{b3:def:spectrum-zariski}
设 \(R\) 为交换含幺环。\(R\) 的\term[sulixiang-pu]{素理想谱}[prime spectrum]是集合
\[
\operatorname{Spec}R=\{\,\mathfrak p\subset R:\mathfrak p\text{ 是素理想}\,\}.
\]
对 \(R\) 的理想 \(I\)，记 \(V(I)=\{\,\mathfrak p\in\operatorname{Spec}R\mid I\subseteq\mathfrak p\,\}\)。以这些 \(V(I)\) 为闭集得到的拓扑称为\term[Zariski-tuopu]{Zariski 拓扑}[Zariski topology]。对单个元素或有限个元素，简记 \(V(f)=V\bigl((f)\bigr)\)、\(V(f_1,\ldots,f_r)=V\bigl((f_1,\ldots,f_r)\bigr)\)。对 \(f\in R\)，记
\[
D(f)=\operatorname{Spec}R\setminus V(f)
=\{\,\mathfrak p:f\notin\mathfrak p\,\},
\]
称为\term[jiben-kaiji]{基本开集}[basic open subset]，也称主开集\termidx[zhukaiji]{主开集}。
\end{definition}
\symidx{\(\operatorname{Spec}R\)：交换环的素理想谱}
\symidx{\(V(I)\)：包含理想 \(I\) 的素理想组成的闭集}
\symidx{\(D(f)\)：素谱中的基本开集}

这里的规则确实定义拓扑，因为
\begin{equation}\label{b3:eq:zariski-closed-identities}
\begin{gathered}
V(0)=\operatorname{Spec}R,\qquad V(R)=\varnothing,
\\
\bigcap_\lambda V(I_\lambda)=V\left(\sum_\lambda I_\lambda\right),
\qquad V(I)\cup V(J)=V(IJ).
\end{gathered}
\end{equation}
前三式来自包含关系。最后一式是因为素理想包含 \(IJ\) 当且仅当它包含 \(I\) 或 \(J\)：若两者都不包含，分别从 \(I,J\) 中选取不属于该素理想的元素，它们的乘积属于 \(IJ\)，与素性矛盾。任意理想之和中的每个元素仅为有限和，所以这也说明无限交闭集的条件成立。

\ProseParagraphBreak
由根的交表示，还有
\begin{equation}\label{b3:eq:zariski-radical-order}
V(I)\subseteq V(J)\quad\Longleftrightarrow\quad
\sqrt J\subseteq\sqrt I.
\end{equation}
因此 \(V(I)=V(\sqrt I)\)，根理想与闭子集反向一一对应。另一方面，\(D(f)\cap D(g)=D(fg)\)，且
\(\operatorname{Spec}R\setminus V(I)=\bigcup_{f\in I}D(f)\)，故基本开集构成开集的一组基。特别地，\(D(f)=\varnothing\) 当且仅当 \(f\) 幂零，\(D(f)=\operatorname{Spec}R\) 当且仅当 \(f\) 是单位。

\subsection{拟紧性}
\label{b3:subsec:0103:02}

设 \(X\) 为拓扑空间。若 \(X\) 的每个开覆盖都有有限子覆盖，则称 \(X\) 为\term[nijin-kongjian]{拟紧空间}[quasi-compact space]。本书不在此词中附加 Hausdorff 条件。

\begin{proposition}\label{b3:prop:spec-quasicompact}
设 \(R\) 为交换含幺环，\(f\in R\)。则 \(\operatorname{Spec}R\) 及其基本开集 \(D(f)\) 都拟紧。
\end{proposition}
\begin{proof}
先将任一开覆盖细化为基本开覆盖。若 \(\operatorname{Spec}R=\bigcup_\lambda D(f_\lambda)\)，则不存在包含全部 \(f_\lambda\) 的素理想。由\cref{b3:thm:maximal-existence}，理想 \((f_\lambda)_\lambda\) 必为 \(R\)，所以一是其中有限个元素的组合。这有限个基本开集已经覆盖全谱，再取它们所属的原开集便得有限子覆盖。

对 \(D(f)\) 的任意子空间开覆盖，可先细化成由基本开集与 \(D(f)\) 的交组成的覆盖；交仍是基本开集，因为 \(D(f)\cap D(g)=D(fg)\)。因此只需考虑 \(D(f)\subseteq\bigcup_\lambda D(g_\lambda)\)。补集的包含关系是 \(V\bigl((g_\lambda)_\lambda\bigr)\subseteq V(f)\)。由\eqref{b3:eq:zariski-radical-order}，某个 \(f^N\) 属于 \((g_\lambda)_\lambda\)，于是属于其中有限个 \(g_\lambda\) 生成的理想。取基本开集又得到有限子覆盖。\(D(f)\) 为空时同样成立。
\end{proof}

\subsection{不可约闭集、一般点与闭点}
\label{b3:subsec:0103:03}

基本开覆盖可以由有限条代数关系控制。闭集还有另一种结构问题：它是否由几个更小的闭集拼成？若不能，下面的判据会把这个闭集辨认为某个素理想的闭包。

\begin{definition}\label{b3:def:irreducible-space}
设 \(X\) 为非空拓扑空间。若 \(X\) 不能写成两个真闭子集之并，则称 \(X\) 为\term[bukeyue-kongjian]{不可约空间}[irreducible space]。闭子集的不可约性用其子空间拓扑理解。若 \(X\) 不能分成两个非空、不交的开集，则称 \(X\) 为连通；不可约性比连通性强，因为在不可约空间中任意两个非空开集必须相交。
\end{definition}

\begin{proposition}\label{b3:prop:spec-irreducible}
设 \(R\) 为交换含幺环，\(I\) 为 \(R\) 的理想。若闭集 \(V(I)\subseteq\operatorname{Spec}R\) 非空，则 \(V(I)\) 不可约当且仅当 \(\sqrt I\) 为素理想。
\end{proposition}
\begin{proof}
若 \(\sqrt I=\mathfrak p\) 素，且 \(V(I)=\bigl(V(I)\cap V(J)\bigr)\cup\bigl(V(I)\cap V(K)\bigr)\)，则 \(V(\mathfrak p)\subseteq V(JK)\)。由\eqref{b3:eq:zariski-radical-order}，\(JK\subseteq\mathfrak p\)，故 \(J\subseteq\mathfrak p\) 或 \(K\subseteq\mathfrak p\)。相应的一个闭子集就是整个 \(V(I)\)。

反过来，若 \(\sqrt I\) 不素，选 \(a,b\notin\sqrt I\)、\(ab\in\sqrt I\)。有 \(V(I)\subseteq V(ab)=V(a)\cup V(b)\)，但 \(V(I)\) 不包含于任何一边，否则根的交表示会使 \(a\) 或 \(b\) 属于 \(\sqrt I\)。这给出两个真闭子集之并，违反不可约性。
\end{proof}

设 \(X\) 为拓扑空间，\(x\in X\)。若单点集 \(\{x\}\) 在 \(X\) 中为闭集，则称 \(x\) 为\term[bidian]{闭点}[closed point]。素谱中还常有闭包比单点大的点；这种闭包可以直接由理想的包含关系读出。

\begin{proposition}\label{b3:prop:generic-and-closed-points}
设 \(R\) 为交换含幺环，\(\mathfrak p\) 为 \(R\) 的素理想。它在 \(\operatorname{Spec}R\) 中所对应点的闭包为 \(V(\mathfrak p)\)。因此它是闭点当且仅当 \(\mathfrak p\) 极大；每个不可约闭子集恰有一个闭包为它的点。
\end{proposition}
\begin{proof}
闭集 \(V(I)\) 包含点 \(\mathfrak p\)，当且仅当 \(I\subseteq\mathfrak p\)，这又保证 \(V(\mathfrak p)\subseteq V(I)\)。所以 \(V(\mathfrak p)\) 是包含该点的最小闭集。它只有该点，当且仅当没有更大的真素理想；利用极大理想存在性，这等价于 \(\mathfrak p\) 极大。若闭集不可约，\cref{b3:prop:spec-irreducible}把它唯一写成 \(V(\mathfrak p)\)，根理想与闭集的一一对应保证此 \(\mathfrak p\) 唯一。
\end{proof}

设 \(X\) 为拓扑空间，\(Z\subseteq X\) 为不可约闭子集。若点 \(\eta\in Z\) 满足 \(\overline{\{\eta\}}=Z\)，则称 \(\eta\) 为 \(Z\) 的\term[yiban-dian]{一般点}[generic point]。一般点是相对于指定闭集而言的，并不与闭点互斥：闭点也是其单点闭子集的一般点。例如整环的素谱有一般点 \((0)\)；在 \(\operatorname{Spec}\mathbb Z\) 中，它的闭包是全谱，所有 \((p)\) 则是闭点。不同点具有不同闭包，所以素谱满足 \(T_0\) 分离性质，却一般不满足每点都闭，更不必是 Hausdorff 空间。

\ProseParagraphBreak
设 \(R\) 为交换含幺环，\(\mathfrak p,\mathfrak q\) 为 \(R\) 的素理想。若 \(\mathfrak p\subseteq\mathfrak q\)，则称 \(\mathfrak q\) 是 \(\mathfrak p\) 的\term[teshuhua]{特殊化}[specialization]，并称 \(\mathfrak p\) 是 \(\mathfrak q\) 的\term[yibanhua]{一般化}[generalization]。由闭包公式，这等价于 \(\mathfrak q\in\overline{\{\mathfrak p\}}\)。素理想 \(\mathfrak p\) 作为 \(V(\mathfrak p)\) 的一般点，记录这个不可约闭子集共同满足的代数关系；但这个闭集未必在全谱中极大。例如域 \(k\) 上的 \(\operatorname{Spec}k[x,y]\) 本身不可约，而 \((x)\) 是其真闭子集 \(V(x)\) 的一般点。

\ProseParagraphBreak
一个只有两个点的谱已经能显示这种拓扑与通常坐标点的差别。设 \(k\) 为域，令 \(R=k[t]_{(t)}\)。任一非零元素都可写为 \(t^n u\)，其中 \(n\ge0\)、\(u\in R\) 为单位；这是把分子中的 \(t\) 因子提取出来，而分子余下的因子及分母都不在 \((t)\) 中。若该元素属于素理想，则 \(n\ge1\)，并由素性得该理想含 \(t\)。另一方面，取值同态 \(f/g\mapsto\dfrac{f(0)}{g(0)}\) 将 \(R\) 满射到 \(k\)，核为 \((t)R\)，所以 \((t)R\) 极大。\(R\) 是整环，故 \((0)\) 也是素理想。这就得到\cref{b3:fig:two-point-specialization}中的全部两点。

\begin{figure}[htbp]
\centering
\begin{tikzcd}[column sep=huge]
\eta=(0) \arrow[r,"\text{特殊化}"] & s=(t)R
\end{tikzcd}
\caption{两点素谱的特殊化关系}
\label{b3:fig:two-point-specialization}
\end{figure}

图中的箭头表示 \(s\) 属于 \(\eta\) 的闭包，不表示距离或路径。具体地，
\[
\overline{\{\eta\}}=\{\eta,s\},\qquad
\overline{\{s\}}=\{s\}.
\]
所以三个开集恰为 \(\varnothing\)、\(\{\eta\}\) 和 \(\{\eta,s\}\)。其中 \(\{\eta\}=D(t)\) 非空，却不含全谱的唯一闭点 \(s\)。在一般素谱中，闭点的剩余域也可能比基域大；第六、七章将在代数闭基域的有限型情形把闭点解释为通常的坐标点。

\subsection{Noether 素谱与有限不可约分量}
\label{b3:subsec:0103:04}

不可约闭子集可能严格包含在另一个不可约闭子集中。例如上面的 \(V(x)\) 就包含在不可约的 \(\operatorname{Spec}k[x,y]\) 内。设 \(X\) 为拓扑空间，若不可约闭子集 \(Z\subseteq X\) 不严格包含于任何更大的不可约闭子集，则称 \(Z\) 为 \(X\) 的\term[bukeyue-fenliang]{不可约分量}[irreducible component]。

\ProseParagraphBreak
设 \(R\) 为交换含幺环，若素理想 \(\mathfrak p\) 不严格包含 \(R\) 的任何其他素理想，则称 \(\mathfrak p\) 为\term[jixiao-sulixiang]{极小素理想}[minimal prime ideal]。由\cref{b3:prop:spec-irreducible}与\eqref{b3:eq:zariski-radical-order}，\(V(\mathfrak p)\) 是全谱的不可约分量，当且仅当 \(\mathfrak p\) 是极小素理想：任一更大的不可约闭集都形如 \(V(\mathfrak q)\)，其包含关系恰对应 \(\mathfrak q\subseteq\mathfrak p\)。因此所有素理想都给出各自闭包的一般点，只有极小素理想给出全谱不可约分量的一般点。

\ProseParagraphBreak
不可约分量是否只有有限个，则是另一个问题。设 \(X\) 为拓扑空间，若 \(X\) 的每条闭集降链最终稳定，则称 \(X\) 为\term[Noether-tuopu-kongjian]{Noether 拓扑空间}[Noetherian topological space]。在素谱中，闭集降链对应根理想升链；环的 Noether 条件使这些升链稳定，因而保证素谱是 Noether 拓扑空间。

\begin{proposition}\label{b3:prop:noether-spectrum-components}
设 \(R\) 为交换含幺 Noether 环。则 \(\operatorname{Spec}R\) 是 Noether 拓扑空间，并且是有限多个不可约分量之并。它们恰为 \(V(\mathfrak p_1),\ldots,V(\mathfrak p_r)\)，其中 \(\mathfrak p_i\) 是 \(R\) 的全部极小素理想；每个素理想都包含其中至少一个。零环对应空分解。
\end{proposition}
\begin{proof}
若 \(R=0\)，素谱为空，结论成立。以下设 \(R\ne0\)。

闭集降链由\eqref{b3:eq:zariski-radical-order}对应根理想升链，Noether 升链条件使其稳定。若存在不能表示为有限个不可约闭集之并的闭集，在这些闭集中取一个关于包含关系的极小元 \(Y\)。降链条件保证极小元存在，否则从任一成员出发连续选择更小者，便得到无限严格降链。于是 \(Y\) 的每个真闭子集都有有限不可约分解。\(Y\) 非空且不可能不可约，否则它自己就是一项分解，故可写成两个真闭子集 \(Y_1,Y_2\) 的并。分别分解 \(Y_1,Y_2\)，再将所得有限个不可约闭集合并，就得到 \(Y\) 的分解，矛盾。

因此全谱有有限不可约闭分解。删去重复成员和被其他成员包含者，得到两两不包含的 \(V(\mathfrak p_1),\ldots,V(\mathfrak p_r)\)，其中各 \(\mathfrak p_i\) 由\cref{b3:prop:spec-irreducible}为素理想。若非空不可约闭集 \(Z\) 被闭集 \(F_1,\ldots,F_r\) 覆盖，则
\[
Z=\bigcup_{i=1}^r(Z\cap F_i).
\]
对 \(r\) 归纳可知某个 \(Z\cap F_i=Z\)：\(r=1\) 时显然；\(r>1\) 时，将第一项与余下各项的并看作 \(Z\) 的两个闭子集。不可约性使其中一边等于 \(Z\)；若是后者，再对 \(r-1\) 项使用归纳。因此任一不可约闭集都包含于某个 \(V(\mathfrak p_i)\)。若不可约闭集 \(Z\) 包含 \(V(\mathfrak p_i)\)，便有
\(V(\mathfrak p_i)\subseteq Z\subseteq V(\mathfrak p_j)\)
对某个 \(j\) 成立；两两不包含性迫使 \(j=i\)，故 \(Z=V(\mathfrak p_i)\)。这说明所列成员就是全部不可约分量，也就是全部极小素理想对应的闭集。

最后，对任一素理想 \(\mathfrak q\)，不可约闭集 \(V(\mathfrak q)\) 包含于某个 \(V(\mathfrak p_i)\)，即 \(\mathfrak p_i\subseteq\mathfrak q\)。
\end{proof}

不同不可约分量可以相交。例如域 \(k\) 上 \(k[x,y]/(xy)\) 的两个极小素理想为 \((\bar x)\) 和 \((\bar y)\)，因为包含 \(xy\) 的素理想必含 \(x\) 或 \(y\)，而 \((x)\)、\((y)\) 本身为互不包含的素理想。两条坐标轴分量 \(V(\bar x)\)、\(V(\bar y)\) 在 \((\bar x,\bar y)\) 处相交，所以存在多个不可约分量并不意味着空间不连通。
